\section{Conditional relation ideals and gamma detectors}
\label{sec:gamma-detectors}

The structural theorems above do not decide the arithmetic nature of
\(\gamma\).  They do, however, turn several concrete relations into exact
certificates against algebraicity.  The point of this section is to state
those certificates with their full quantifiers and relation ideals.

Put
\[
 D=\delta,\qquad X=E_1(1),\qquad Y=\Ei(1),\qquad
 Z=\Ci(1),\qquad W=\Chi(1),
\]
and retain the pure exponential field
\(\Eexp=\Qbar(e^\alpha:\alpha\in\Qbar)\).

\begin{proposition}[The conditional ideal over the pure exponential field]
\label{prop:conditional-pure-exp-ideal}
If \(\gamma\in\Qbar\), then \(X,Y,Z\) are algebraically independent over
\(\Eexp\), and the complete relation ideal of \(D,X,Y,Z,W\) over
\(\Eexp\) is
\begin{equation}\label{eq:conditional-pure-exp-ideal}
 (\,\mathsf D-e\mathsf X,\;2\mathsf W-\mathsf Y+\mathsf X\,).
\end{equation}
Consequently
\begin{equation}\label{eq:gamma-or-delta-global}
 \boxed{\ \gamma\text{ is transcendental, or }
 \delta\text{ is transcendental over }\Eexp.\ }
\end{equation}
\end{proposition}

\begin{proof}
Under the hypothesis, the connection identities give
\[
 X=\Ein(1)-\gamma,\qquad
 Y=\gamma-\Ein(-1),\qquad
 Z=\gamma-\frac{\Ein(i)+\Ein(-i)}2.
\]
They are three linearly independent combinations of the four
\(\Ein\)-coordinates at \(1,-1,i,-i\).  Theorem
\ref{thm:pure-exp-base} and Lemma~\ref{lem:linear-image-ein} therefore make
them algebraically independent over \(\Eexp\).  The remaining coordinates
satisfy, and are determined by,
\[
 D=eX,\qquad 2W=Y-X.
\]
This proves the ideal statement.  In particular \(D=eX\) is transcendental
over \(\Eexp\) whenever \(\gamma\) is algebraic, proving
\eqref{eq:gamma-or-delta-global}.
\end{proof}

For \(m\ge1\), set
\[
 T_m=e^{\gamma^m}.
\]

\begin{theorem}[Exact \(T_m\)-ideals under algebraic \(\gamma\)]
\label{thm:Tm-exact-ideals}
Assume \(\gamma\in\Qbar\), fix \(m\ge1\), and put
\(\eta=\gamma^m\).  If \(\eta\notin\Q\), then
\[
 T_m,D,X,Y,Z
\]
are algebraically independent over \(\Qbar\).  If
\(\eta=a/b\in\Q_{>0}\), where \((a,b)=1\), then the complete relation
ideal of these five coordinates is the principal prime ideal
\begin{equation}\label{eq:Tm-rational-ideal}
 (\,\mathsf D^a-\mathsf T^b\mathsf X^a\,).
\end{equation}
After adjoining \(W\), in the first case one adds only
\(2\mathsf W-\mathsf Y+\mathsf X\); in the second case the full ideal is
generated by that linear form and \eqref{eq:Tm-rational-ideal}.

In either case \(T_m\) and \(D\) are algebraically independent.  Hence,
unconditionally,
\begin{equation}\label{eq:Tm-delta-detector}
 \boxed{\ P(e^{\gamma^m},\delta)=0\text{ for some }
 0\ne P\in\Qbar[U,V]\quad\Longrightarrow\quad
 \gamma\text{ is transcendental}.\ }
\end{equation}
\end{theorem}

\begin{proof}
Since \(0<\gamma<1\), the number \(\eta\) is nonzero and differs from
\(1,-1,i,-i\).  Apply Theorem~\ref{thm:master} to
\(\{1,-1,i,-i,\eta\}\).  With
\[
 A=\Ein(1),\qquad B=\Ein(-1),\qquad
 C=\frac{\Ein(i)+\Ein(-i)}2,
\]
the relevant coordinates are birationally related by
\[
 e=\frac D X,\qquad A=X+\gamma,\qquad
 B=\gamma-Y,\qquad C=\gamma-Z.
\]
If \(1\) and \(\eta\) are \(\Q\)-independent, the exponential coordinates
\(e,T_m\) have no toric relation, and the five displayed coordinates are
independent.  If \(\eta=a/b\), their only toric relation is
\(T_m^b=e^a\), which becomes
\(D^a=T_m^bX^a\).  The exponent vector is primitive, so the binomial is
prime; the transcendence-degree formula gives height one and excludes any
additional relation.  The identity \(2W=Y-X\) handles \(W\).

In both cases contraction to \(\Qbar[\mathsf T,\mathsf D]\) is zero.  Thus
a relation as in \eqref{eq:Tm-delta-detector} contradicts the temporary
hypothesis \(\gamma\in\Qbar\).
\end{proof}

An exact zero is not the only possible certificate.  A quantitative
\(E\)-function theorem also rules out fixed-degree super-Liouville
approximation under the algebraicity hypothesis.

\begin{corollary}[Fixed-degree approximation detector]
\label{cor:fixed-degree-gamma-detector}
Fix \(m,d\ge1\).  Suppose that nonzero polynomials
\(P_n\in\Z[U,V]\), of degree at most \(d\), satisfy
\[
 H(P_n)\longrightarrow\infty,\qquad
 P_n(T_m,D)\ne0,
\]
and
\begin{equation}\label{eq:fixed-degree-super-small}
 \frac{-\log|P_n(T_m,D)|}{\log H(P_n)}\longrightarrow\infty.
\end{equation}
Then \(\gamma\) is transcendental.
\end{corollary}

\begin{proof}
Assume \(\gamma\in\Qbar\).  Then
\[
 f_1(z)=e^{\gamma^m z},\qquad
 f_2(z)=e^z\bigl(\Ein(z)-\gamma\bigr)
\]
are \(E\)-functions with algebraic coefficients and
\((f_1(1),f_2(1))=(T_m,D)\).  Theorem~\ref{thm:Tm-exact-ideals} makes the
two values algebraically independent.  The fixed-degree case of the
Adamczewski--Faverjon algebraic-independence measure
\cite{AdamczewskiFaverjon2025} supplies constants \(c_d,\kappa_d>0\) with
\[
 |P(T_m,D)|\ge c_dH(P)^{-\kappa_d}
\]
for every nonzero integral polynomial of degree at most \(d\).  This
contradicts \eqref{eq:fixed-degree-super-small}.
\end{proof}

\begin{remark}[Logical boundary]
Theorems~\ref{thm:Tm-exact-ideals} and
Corollary~\ref{cor:fixed-degree-gamma-detector} are one-way detectors.
They do not construct a polynomial relation or a super-small sequence.
Moreover \(e^{\gamma^m}\) is an \(E\)-value with algebraic coefficients
only inside the contradiction hypothesis.  Treating it as an unconditional
\(E\)-value would reverse the logic.
\end{remark}
